\chapter{Канонический шумовой кадр общего KMS-следа}
\label{ch:v8-canonical-noise-frame-common-trace-gate}

\section{От списка скачков к шумовому модулю}

Предыдущий гейт показал, что коэффициенты зависят от нормировки операторов
скачка. Вместо покомпонентного соглашения рассмотрим весь уже полученный
шумовой подмодуль и метрику, индуцированную верным KMS-состоянием
$\rho_N$:
\begin{equation}
 G_{ab}=\Tr(\rho_N F_a^*F_b).
 \label{eq:v8-common-trace-noise-gram}
\end{equation}
Ортонормированный кадр этого Грам-оператора определён с точностью до
унитарного преобразования. Сумма GKSL-диссипаторов по такому кадру от этого
преобразования не зависит и потому задаёт канонический Казимир шумового
подпространства.

Исходный список содержит 25 компонент: тринадцать компонент переноса и двенадцать
калибровочных компонент. Однако грамова матрица переноса имеет ранг семь и нульность шесть.
Шесть нулей точно соответствуют парам вещественной и мнимой вариаций одного
комплексного оператора: $V$ и $iV$ порождают один комплексный шумовой канал,
а их раздельный подсчёт удваивал его вес.

После факторизации получаем
\begin{equation}
 \dim_{\mathbb C}\mathcal N_{\rm tr}=7,
 \qquad
 \dim_{\mathbb C}\mathcal N_{\rm gauge}=12,
 \qquad
 \dim_{\mathbb C}\mathcal N=19.
 \label{eq:v8-common-trace-noise-dimension}
\end{equation}
Остатки ортонормировки равны $6.18\cdot10^{-16}$ и
$9.54\cdot10^{-16}$ соответственно.

\section{Базис-независимый генератор}

Пусть $\{W_a\}_{a=1}^{7}$ и $\{T_b\}_{b=1}^{12}$ --- ортонормированные в
общей KMS-метрике реперы переноса и калибровочные реперы. Определим
\begin{equation}
 \mathcal L_{\rm tr}=
 \sum_{a=1}^{7}\mathcal D^{\rm KMS}_{W_a},
 \qquad
 \mathcal L_{\rm gauge}=\sum_{b=1}^{12}\mathcal D_{T_b},
 \qquad
 \mathcal L_{\rm can}=\mathcal L_{\rm tr}+\mathcal L_{\rm gauge}.
 \label{eq:v8-canonical-noise-casimir}
\end{equation}
В части переноса сохранено уже выведенное отношение прямой и обратной
скорости $e^{-2}$.

Четыре случайные невырожденные замены исходного кадра изменили
$\mathcal L_{\rm can}$ не более чем на $5.34\cdot10^{-16}$. Таким образом,
произвольное масштабирование и смешивание координат действительно удалены.

Полный генератор унитален, KMS-симметричен и примитивен:
\begin{align}
 \dim\ker\mathcal L_{\rm can}&=1,\\
 \operatorname{gap}(-\mathcal L_{\rm can})&=0.0207491297,\\
 \|G\mathcal L_{\rm can}-\mathcal L_{\rm can}^*G\|
 &=3.40\cdot10^{-17}.
 \label{eq:v8-canonical-noise-generator-diagnostics}
\end{align}

\section{Выбранная точка шестимерного конуса}

Канонический Казимир с машинным относительным остатком
$5.89\cdot10^{-16}$ лежит в прежнем шестисемейном конусе. В порядке
$(L,SU(3),SU(2),U(1),QLYR,XLdR)$ он выбирает ненормированные веса
\begin{equation}
 (0.83698439,\ 10.88079708,\ 4.81111832,\ 4.52410355,\
 5.44039854,\ 5.44039854).
 \label{eq:v8-canonical-trace-family-weights}
\end{equation}
После нормировки их сумма равна единице:
\begin{equation}
 (0.0262100,\ 0.3407298,\ 0.1506591,\ 0.1416713,\
 0.1703649,\ 0.1703649).
 \label{eq:v8-canonical-trace-normalized-weights}
\end{equation}
Равенство последних двух компонент теперь является следствием общего
следа, а не ручным соглашением. Все пять относительных свобод предыдущего
конуса устранены; остаётся только общий множитель физического времени.

\section{Граница физического статуса}

Математическая каноничность не равна физической необходимости. Общий
KMS-след задаёт естественную изотропную метрику на уже выбранном шумовом
подпространстве, но фактическая калибровочная симметрия сохраняет его
разложение на неэквивалентные компоненты и сама по себе допускает разные
положительные веса этих компонент. Следовательно, получен сильный
канонический представитель, но ещё не теорема о том, что природа обязана
выбрать именно изотропный по следу Казимир.

\begin{theorem}[Канонический оператор Казимира шумового факторпространства]
Общий KMS-след на уже допущенном связывающе-калибровочно-межсекторном шумовом подмодуле
после факторизации по шести комплексно-линейным дублям задаёт
базис-независимый 19-мерный ортонормированный кадр. Его Казимир порождает
примитивный KMS-симметричный генератор и единственную точку прежнего
шестимерного конуса с точностью до общего масштаба времени.
\end{theorem}

\begin{nogocriterion}[Запрет объявления изотропии по следу физической симметрией]
Нельзя считать веса \eqref{eq:v8-canonical-trace-family-weights}
единственно физическими, пока не показано, что допустимая симметрия смешивает
весь 19-мерный шумовой модуль неприводимо или что одно родительское действие
непосредственно порождает его ортогональный проектор. Общий след даёт
канонический минимум предпосылок, но не заменяет такую теорему.
\end{nogocriterion}

\begin{protocol}
\begin{enumerate}
 \item сырых шумовых компонент: \textbf{25};
 \item нулевых комбинаций переноса: \textbf{6};
 \item физическое комплексное шумовое факторпространство: \textbf{19};
 \item общая KMS-грамова форма: \textbf{положительна на факторпространстве};
 \item остаток ортонормировки: \textbf{меньше $10^{-15}$};
 \item инвариантность при смене кадра: \textbf{пройдена};
 \item неподвижная алгебра: \textbf{$\mathbb C I_{21}$};
 \item щель в следовых единицах времени: \textbf{$0.0207491297$};
 \item пять относительных весов: \textbf{канонически зафиксированы};
 \item абсолютный физический масштаб времени: \textbf{открыт};
 \item физическая обязательность изотропии по следу: \textbf{не доказана};
 \item следующий гейт: \textbf{симметрийный допуск изотропии шумового модуля}.
\end{enumerate}
\end{protocol}

Машинный сертификат сохранён в
\path{s2t_v8_canonical_noise_frame_common_trace_gate_results.json}.

\begin{thebibliography}{9}
\bibitem{v8-noise-frame-gks}
V.~Gorini, A.~Kossakowski, E.~C.~G.~Sudarshan,
\emph{Completely positive dynamical semigroups of $N$-level systems},
J. Math. Phys. 17 (1976) 821--825.
\bibitem{v8-noise-frame-lindblad}
G.~Lindblad,
\emph{On the generators of quantum dynamical semigroups},
Commun. Math. Phys. 48 (1976) 119--130.
\end{thebibliography}